\documentclass[11pt]{article}
\usepackage[T1]{fontenc}
\usepackage[utf8]{inputenc}
\usepackage{lmodern,amsmath,amssymb,amsthm,mathtools,esint}
\usepackage[margin=1in]{geometry}
\usepackage[hidelinks]{hyperref}
\providecommand{\tightlist}{\setlength{\itemsep}{0pt}\setlength{\parskip}{0pt}}
\setlength{\emergencystretch}{2em}
\title{Two Proof Increments for the Weak Endpoint and Extremal Questions}
\author{Qiushi Line 1: research supplement}
\date{29 September 2026}
\begin{document}
\maketitle
\noindent\textbf{Scope.} This supplement proves two intermediate upgrades from explicitly identified estimates in the supplied integrated manuscript. It does not assert a solution of unrestricted O10 or the optimal coefficient in O11. The manuscript's inherited inputs are not independently re-proved here.
\tableofcontents
\bigskip
\hypertarget{o10-a-strong-current-upgrade-from-the-actual-constant-weight-equation}{%
\section{O10: a strong-current upgrade from the actual constant-weight
equation}\label{o10-a-strong-current-upgrade-from-the-actual-constant-weight-equation}}

Date: 2026-09-29. Source:
the supplied integrated manuscript (source hash recorded in the accompanying status file), especially
\texttt{sec:single-parameter}, \texttt{sec:endpoint-covariance},
\texttt{sec:phase-curvature}, \texttt{sec:torsion-prolongation}, and
\texttt{sec:quartic-flux}.

\hypertarget{result-and-scope}{%
\subsection{Result and scope}\label{result-and-scope}}

There is a new global estimate available before attempting any further
manipulation of the final quartic flux identity. The actual
constant-weight positive spinor has uniformly integrable fourth-power
first derivative along every hypothetical O10 endpoint sequence. This
upgrades the cubic torsion Dirac equation, its length, the original
current derivative, and the full first torsion projection to strong
critical statements. It uses no diameter, Sobolev, noncollapse,
pointwise Ricci bound, or Euler-density upper bound.

It does \textbf{not} complete O10: it does not estimate the fourth-power
derivative of the cubic torsion, its relative-index flux, or the
remaining curvature-weighted radial-horizontal interaction. In
particular it does not contradict the shrinking actual local torsion
cores already constructed in the manuscript.

All averages below use the actual total volume, even when the manifold
varies. Thus every vanishing squared norm is an \texttt{o(V)} statement,
not an unnormalized smallness assumption.

\hypertarget{a-slow-gradient-truncation-lemma}{%
\subsection{1. A slow gradient truncation
lemma}\label{a-slow-gradient-truncation-lemma}}

Let \((M_j,g_j)\) be closed Riemannian manifolds of fixed dimension and
\(E_j\) Hermitian bundles of fixed rank, with unitary connections. Let
\(u_j\) satisfy

\[
 |u_j|\le C_0,\qquad \nabla^*\nabla u_j=f_j(|u_j|^2)u_j,
 \qquad |f_j|+|f'_j|\le C_0.
\]

Suppose

\[
 \operatorname{Ric}_{g_j}=R_{0,j}+e_j,
 \quad F_j=F_{0,j}+E_j,
 \quad |R_{0,j}|+|F_{0,j}|\le C_0,
\]

where the symbol \(E_j\) on the right denotes a curvature error, and
assume

\[
 \mu_j:=\fint(|e_j|^2+|E_j|^2)\longrightarrow0,
 \qquad \sup_j\mu_j\le C_0.
\]

No independent Hessian or fourth-power gradient budget is needed: these
follow from the displayed hypotheses by the unweighted version of the
argument below.

There is \(\epsilon>0\), depending only on the displayed constants,
dimension, and rank, such that, with

\[
 q=|\nabla u|^2,\quad H=\nabla^2u,\quad
 w_L=\left(\frac q{L^2+q}\right)^\epsilon,
 \qquad L\ge1,
\]

one has

\[
 \boxed{\fint w_L(|H|^2+q^2)
              \le C L^{-2\epsilon}+C\mu_j.}\tag{1}
\]

Consequently

\[
 \boxed{\lim_{L\to\infty}\limsup_j
       \fint_{\{|\nabla u_j|\ge L\}}|\nabla u_j|^4=0.}\tag{2}
\]

This establishes uniform integrability of \(|\nabla u_j|^4\) relative to
normalized volume. Formula (1) also controls \(|H|^2\) on the
large-gradient set. It does \textbf{not}, by itself, prove uniform
integrability of \(|H|^2\) on arbitrary small-volume sets.

\hypertarget{proof}{%
\subsubsection{Proof}\label{proof}}

Use the smooth regularization

\[
 w_{L,\delta}=\left(\frac{q+\delta}{L^2+q+\delta}\right)^\epsilon,
 \quad \delta>0.
\]

It obeys

\[
 |dw_{L,\delta}|\,|\nabla u|\le2\epsilon w_{L,\delta}|H|.
\tag{3}
\]

Put \(P_i=\nabla_i u\), \(H_{ij}=\nabla_iP_j\), and \(I=\int w|H|^2\),
\(J=\int wq^2\), \(A_w=\int w(1+q)\). Antisymmetry gives the exact
algebraic identity

\[
 I=\int w\operatorname{Re}\langle H_{ij},H_{ji}\rangle
       +\frac12\int w\sum_{i,j}|F_{ij}u|^2.
\tag{4}
\]

In the first term integrate the derivative indexed by \(j\) in
\(H_{ji}\). Then commute the two derivatives in
\(\nabla_jH_{ij}=\nabla_j\nabla_iP_j\). The commutator on the
covector-valued section \(P\) contains only
\(\operatorname{Ric}*P+F*P\), and \(\sum_j\nabla_jP_j=-f(|u|^2)u\). In
particular \textbf{no derivative of either curvature occurs}. The
derivative of \(f(|u|^2)u\) is bounded by \(C|P|\). The derivative of
the scalar weight contributes at most \(2\epsilon I\), by (3). Thus,
after taking \(\epsilon\) sufficiently small,

\[
 I\le C A_w+C\int w(|e|+|E|)q+C\int (|e|^2+|E|^2).
\tag{5}
\]

The term involving \(F^2\) in (4) was bounded using \(F=F_0+E\) and
bounded \(F_0\). No derivative of the bounded model is required.

For completeness, a weighted bounded-section interpolation follows by
integrating the divergence of \(wq\operatorname{Re}\langle u,P\rangle\).
The derivative of \(q\) contributes at most \(C\int wq|H|\); the
derivative of \(w\) has the same bound by (3); the trace Hessian term is
bounded by \(C\int wq\). Therefore

\[
 J\le C\int wq+C\sqrt{IJ},\qquad
 J\le C_1(I+A_w).\tag{6}
\]

Since \(0\le w\le1\), for any fixed \(a>0\),

\[
 \int w(|e|+|E|)q
 \le aJ+C_a\int(|e|^2+|E|^2).
\tag{7}
\]

Choose \(a\) after \(C_1\) so that (5)--(7) absorb the resulting
multiple of \(I\). It follows that

\[
 I+J\le C A_w+C\mu_jV_j.
\tag{8}
\]

Before estimating the weighted right side, take \(w\equiv1\) in
(4)--(8). Testing the section equation by \(u\) gives \[
\int q=\int f(|u|^2)|u|^2\le CV.
\] The same absorption then proves \[
\fint(|H|^2+q^2)\le C(1+\mu_j)\le C.
\] Thus the energy budget used in the following step is a consequence,
not an extra hypothesis.

For \(0<\epsilon\le1\),

\[
 w_{L,\delta}\le L^{-2\epsilon}(q+\delta)^\epsilon.
\]

The just-proved normalized \(q^2\) bound controls the average of
\((q+\delta)^\epsilon(1+q)\), uniformly for \(0<\delta\le1\). Divide (8)
by \(V_j\), then let \(\delta\downarrow0\), obtaining (1). On
\(q\ge L^2\), \(w_L\ge2^{-\epsilon}\), so (2) follows. Every calculation
is on a closed manifold with a globally smooth regularized scalar
weight.

\hypertarget{apply-the-lemma-to-the-hypothetical-o10-endpoint}{%
\subsection{2. Apply the lemma to the hypothetical O10
endpoint}\label{apply-the-lemma-to-the-hypothetical-o10-endpoint}}

Use the manuscript's exact Yamabe reduction and a single ordinary
solution rescaled to

\[
 s=-4,\quad D_A\Phi=0,\quad F_A^+=2i\sigma(\Phi),
 \quad \rho=|\Phi|^2\le2.
\]

Set \(P=\nabla_A\Phi\), \(H=\nabla_AP\), \(\zeta=\langle\Phi,P\rangle\),
\(\tau=t_\Phi\), \(m=|\tau|^2\), \(\rho_*=4/3\), \(m_*=128/81\), and
\(\eta=r+D_\beta/V+\beta^{-2}\to0\).

The actual semilinear equation is

\[
 \nabla_A^*\nabla_A\Phi=(1-\rho/2)\Phi.\tag{9}
\]

The latest bounded ASD fit supplies the lemma's full
connection-curvature hypothesis. The old Hessian budget is consistent
with, and can also be recovered from, the unweighted part of the lemma.
With

\[
 f=-iF_A^-,\quad Q=-iq_-(\tau),\quad
 S=m_*+2\rho^2|P|^2,\quad
 K=-\frac98\frac{\rho^3Q}{S},
\]

the manuscript proves

\[
 |K|\le C,\qquad \fint|f-K|^2\le C\sqrt\eta.
\tag{10}
\]

The traceless curvature of \(S^+\) contains \(W^++s/12\) and the
off-diagonal Ricci block; it contains \textbf{no \(W^-\)}. Its central
part is \(F_A/2\). Retain the scalar block, the bounded actual self-dual
determinant curvature, and the bounded ASD model \(iK\) as \(F_0\). The
remaining curvature has

\[
 \fint|F^{S^+}-F_0|^2\le C\sqrt\eta,
 \quad |F_0|\le C,
 \quad \fint|\operatorname{Ric}+g|^2\le\eta.
\tag{11}
\]

Here \(W^+\to0\) follows from the same-solution parameter estimate,
including its bounded axial model of size \(O(\beta^{-1})\). Applying
(1) gives

\[
 \boxed{\fint w_L(|H|^2+|P|^4)
         \le C L^{-2\epsilon}+C\sqrt\eta.}\tag{12}
\]

There is no hypothesis on the opposite Weyl curvature and hence no use
of \(\chi/A\le B\) in this result.

\hypertarget{strong-critical-consequences}{%
\subsection{3. Strong critical
consequences}\label{strong-critical-consequences}}

\hypertarget{the-mixed-spinor-product-and-torsion-dirac-residual}{%
\subsubsection{3.1 The mixed spinor product and torsion Dirac
residual}\label{the-mixed-spinor-product-and-torsion-dirac-residual}}

The old identity gives \(\fint|\zeta|^2\le C\eta\), while
\(|\zeta|\le\sqrt2|P|\). Splitting into \(|P|\le L\) and its complement
in (12) yields

\[
 \fint|\zeta|^2|P|^2
 \le C\{L^2\eta+L^{-2\epsilon}+\sqrt\eta\}\to0.
\tag{13}
\]

For example choosing \(L=\eta^{-1/(2+2\epsilon)}\) gives a bound by
\(C\eta^{\epsilon/(1+\epsilon)}+C\sqrt\eta\). Also
\(\fint|\zeta|^4\to0\). The exact cubic equation, with axial
cancellation retained, now implies

\[
 \boxed{\|D^-_{3A}\tau\|_{L^2_{\rm avg}}\to0.}\tag{14}
\]

This is stronger than the previous \(L^q\), \(q<2\), conclusion.

\hypertarget{torsion-length-and-positive-spinor-energy-density}{%
\subsubsection{3.2 Torsion length and positive-spinor energy
density}\label{torsion-length-and-positive-spinor-energy-density}}

The existing torsion-length concentration gives \(m\to m_*\) in
normalized measure. Since \(m\le2\rho^2|P|^2\le8|P|^2\), (12) makes
\(m^2\) uniformly integrable. Vitali on the probability spaces
\((M_j,dV_j/V_j)\) gives

\[
 \boxed{\|m-m_*\|_{L^2_{\rm avg}}\to0.}\tag{15}
\]

The exact algebraic defect satisfies

\[
 \nu=2\rho^2|P|^2-m,\quad
 0\le\nu\le4\rho|\zeta|^2,
\]

so \(\nu\to0\) in normalized \(L^2\). On \(\rho\ge c>0\), use (15) and
\(\rho\to\rho_*\) in measure; on \(\rho<c<\rho_*\), its relative volume
tends to zero and (12) controls \(|P|^4\). Consequently

\[
 \boxed{\big\||P|^2-4/9\big\|_{L^2_{\rm avg}}\to0.}\tag{16}
\]

The scalar equation \(\Delta\rho=-2|P|^2+2\rho-\rho^2\) therefore gives

\[
 \|\Delta\rho\|_{L^2_{\rm avg}}\to0.\tag{17}
\]

\hypertarget{the-full-current-derivative}{%
\subsubsection{3.3 The full current
derivative}\label{the-full-current-derivative}}

The actual polynomial curvature equation is

\[
 \rho^3 f+2Q=2\rho^2(d\gamma)^-+\mathcal R,
 \quad \gamma=\operatorname{Im}\zeta,
 \quad |\mathcal R|\le
 C(\rho^3|\operatorname{Ric}^{\circ}|+
              \rho^{3/2}|\zeta||P|).
\tag{18}
\]

Real/imaginary conventions may reverse both signs of the derivative
term; only its norm is used here. By (13), \(\mathcal R\to0\) in
normalized \(L^2\). Furthermore

\[
 \rho^3 f+2Q
 =\rho^3(f-K)+\left(2-\frac98\frac{\rho^6}{S}\right)Q.
\tag{19}
\]

The first term is small by (10). The coefficient in the second is
uniformly bounded, converges to zero in measure, and \(|Q|^2=m^2/8\) is
uniformly integrable by (15). Thus the second term is small in
normalized \(L^2\). Equation (18) proves \(\rho^2(d\gamma)^-\to0\) in
that norm.

This weighted assertion becomes unweighted without assuming a pointwise
lower bound for \(\rho\). Indeed

\[
 |d\zeta|\le C(|P|^2+\rho|F^{S^+}|).
\tag{20}
\]

The squares on the right are uniformly integrable, respectively by (12)
and the bounded-model fit (11). The set \(\{\rho<c\}\), for any fixed
\(0<c<\rho_*\), has vanishing relative volume. On its complement divide
by \(\rho^2\); on that set use (20). Hence

\[
 \|(d\gamma)^-\|_{L^2_{\rm avg}}\to0.
\tag{21}
\]

The real part \(\operatorname{Re}\zeta=d\rho/2\) is exact, and
\(\gamma\) is a globally smooth one-form with \(\delta\gamma=0\). Stokes
gives exactly

\[
 \int|d\gamma^+|^2=\int|d\gamma^-|^2,
\]

so \(d\gamma\to0\) in normalized \(L^2\). The one-form and scalar
Bochner formulas yield

\[
 \int|\nabla\gamma|^2
 =\int|d\gamma|^2-\int\operatorname{Ric}(\gamma,\gamma),
\]

\[
 \int|\nabla^2\rho|^2
 =\int(\Delta\rho)^2-\int\operatorname{Ric}(d\rho,d\rho).
\]

The \(-g\) terms are small by the original current bound. The
Ricci-error pairings are bounded by its \(L^2_{\rm avg}\) norm times the
squared \(L^4_{\rm avg}\) norm of \(\gamma\) or \(d\rho\), which tend to
zero by (13). Therefore

\[
 \boxed{\|\nabla\zeta\|_{L^2_{\rm avg}}\to0.}\tag{22}
\]

This proves the conclusion sought from the manuscript's current--Hessian
correlation \(\mathcal B\) \textbf{without adding \(\mathcal B\to0\) as
a hypothesis} and without claiming that \(\mathcal B\) itself has been
estimated.

\hypertarget{the-full-first-torsion-projection}{%
\subsubsection{3.4 The full first torsion
projection}\label{the-full-first-torsion-projection}}

Use the already proved actual identity

\[
 N=\sqrt{\rho/2}\,\mathcal L(\nabla\zeta)+\mathcal R_N,
\quad
 |\mathcal R_N|\le C\{\rho^{3/2}(
 |\operatorname{Ric}^{\circ}|+|W^+-W_0|)+|\zeta||P|\}.
\]

Equations (13), (22), and the curvature errors give

\[
 \boxed{\|N\|_{L^2_{\rm avg}}\to0.}\tag{23}
\]

This is the complete displayed first torsion projection, not merely its
Dirac trace.

\hypertarget{the-remaining-torsion-gradient-has-a-fixed-nonzero-mean}{%
\subsubsection{3.5 The remaining torsion gradient has a fixed nonzero
mean}\label{the-remaining-torsion-gradient-has-a-fixed-nonzero-mean}}

The unweighted twisted Lichnerowicz identity is now legitimate at the
limiting strength:

\[
 \fint|\nabla_{3A}\tau|^2
 =\fint|D^-_{3A}\tau|^2+\fint m
       -\frac32\fint\langle c(F_A^-)\tau,\tau\rangle.
\]

Replace \(f=-iF_A^-\) by its bounded fit \(K\) using Cauchy--Schwarz
with the uniformly bounded \(L^2\) norm of \(m\). The Clifford identity
and (15) then give

\[
 \frac32\fint\langle c(F_A^-)\tau,\tau\rangle
 =-\frac{27}{32}\fint\frac{\rho^3m^2}{S}
 \longrightarrow-m_*.
\]

The bounded coefficient \(\rho^3/S\), convergence in measure, and
uniform integrability of \(m^2\) justify this replacement. Consequently

\[
 \boxed{\fint|\nabla_{3A}\tau|^2\longrightarrow
              2m_* =256/81.}\tag{25}
\]

Thus the strengthened first projection does not imply that the torsion
is parallel: its surviving derivative energy is nonzero and sits in the
complementary type component.

\hypertarget{what-has-actually-been-removed-and-what-remains}{%
\subsection{4. What has actually been removed, and what
remains}\label{what-has-actually-been-removed-and-what-remains}}

The previous chain had only subcritical torsion-length convergence and
weak Dirac control, and had left \(\nabla\zeta\) behind an unresolved
current--Hessian covariance. Equations (14)--(16), (22), and (23)
genuinely strengthen that chain.

The critical quantity needed for the final signed characteristic flux is
still of higher differential order. Neither

\[
 \fint|D\tau|^2\to0,\quad
 \fint|N|^2\to0,\quad
 \fint(m-m_*)^2\to0
\]

controls

\[
 \fint_{\{a<m<b\}}
      |\langle\tau,\nabla\tau\rangle|^2|\nabla\tau|^2,
\tag{24}
\]

or the curvature-weighted radial--horizontal product in
\texttt{proof:phase-curvature:e3.6}. The first quantity involves fourth
powers of \textbf{\(\nabla\tau\)}, which contains the original Hessian.
Formula (12) controls fourth powers of \textbf{\(\nabla\Phi\)}; these
must not be conflated.

The local actual simple torsion-zero cores remain compatible with
(14)--(23): their contributions to the strong quadratic norms vanish
under shrinking, while their transition flux stays of order one. Thus
one cannot append an unsupported topological vanishing conclusion to the
new strong-current theorem.

The shortest remaining global step is to control the signed interaction
in the final identity

\[
 \|\vartheta W^--\widetilde H\|_2^2-
 \|\widetilde H\|_2^2
 =\|\vartheta\mathsf R\|_2^2-
 2\int\vartheta\langle\mathsf R,H^+\rangle
\]

using (14)--(23), the full actual commutator/bracket, and the exterior
matching. The old positive \(20V/3\) has already canceled in this
identity and cannot be reinserted. A proof of a uniform signed
separation, or a complete calibrated closed escape construction, is
still required for O10.

\hypertarget{independent-verification-within-this-turn}{%
\subsection{Independent verification within this
turn}\label{independent-verification-within-this-turn}}

The slow-truncation lemma was independently checked by the root agent
and \texttt{o10\_counterexample}. The latter supplied the
derivative-free weighted Hessian interchange (4), which is stronger and
cleaner than commuting the rough Laplacian and integrating a
derivative-of-curvature term. Both checks explicitly preserve the
distinction between gradient fourth-power uniform integrability and full
Hessian-square uniform integrability.

\clearpage
% Standalone insertion fragment. Requires amsmath, amssymb, and amsthm.
% All notation is defined below; no project-specific macros are used.
% This fragment does not assert a solution of the unrestricted coefficient problem.

\section{Variational selection at the topological defect scale}
\label{sec:o11-variational-selection-20260929}

We extract additional variational information from a sufficiently accurate
minimizing sequence. The argument uses finite metric variations directly;
it does not require a smooth minimizer, an Ekeland principle, or geometric
compactness. The test class is critical under four-dimensional scaling and
therefore admits variations supported on shrinking regions.

\paragraph{Setting and inherited estimates.}
Let
\[
 M_j=X_j\mathbin{\#}k_j\overline{\mathbb{CP}}^{\,2}
       \mathbin{\#}\ell_j(S^1\times S^3),
 \qquad A_j=c_1^2(X_j),\quad Q_j=k_j+4\ell_j>0,
 \quad h_j=Q_j/A_j\longrightarrow0,
\]
be the balanced complex-hyperbolic surgery family. For a smooth metric set
\[
 P(g)=\int_{M_j}|\operatorname{Ric}^{\circ}_g|^2\,dV_g,
 \qquad R(g)=\int_{M_j}|\operatorname{Ric}_g|^2\,dV_g,
 \qquad S(g)=\int_{M_j}s_g^2\,dV_g.
\]
Ricci norms are tensor norms. Curvature and Weyl norms below are the
Hilbert--Schmidt norms of their operators on two-forms. Define
\[
 p_j=\inf_{g\text{ on }M_j}P(g),\qquad
 \mathcal E(g)=4-
 \frac{\bigl(\int_{M_j}s_g\,dV_g\bigr)^2}
      {\operatorname{Vol}(g)R(g)},\qquad
 \alpha_j=\inf_{g\text{ on }M_j}\mathcal E(g).
\]
We use the following established properties of this family:
\begin{enumerate}
\item Every conformal class is negative Yamabe. Passing to its negative
Yamabe representative does not increase $P$, since
\[
 P(g)-P(\widehat g)=\frac{S(g)-S(\widehat g)}{12}\ge0.
\]
The representative can be scaled to $s_{\widehat g}=-4$ without changing
$P$.
\item The recovery construction gives $p_j\le8\pi^2Q_j$. For every metric,
\[
 8\pi^2A_j+P(g)\le R(g)
 \le8\pi^2A_jU_j+4P(g),\qquad U_j=1+O(h_j).
\]
\item The variational comparison gives
\[
 \alpha_j=\frac{4p_j}{8\pi^2A_j}+O(h_j^2).
\]
Moreover, $\chi(M_j)=A_j/3+O(Q_j)$ and
\[
 \int_{M_j}|\operatorname{Rm}_g|^2\,dV_g
       =8\pi^2\chi(M_j)+P(g).
\]
\end{enumerate}
All constants in these inherited estimates are uniform along the family.
The new argument below is conditional only on these displayed inputs.

\paragraph{Critical test budget.}
For a smooth symmetric covariant two-tensor $H$, put
\[
 \mathfrak D_g(H)=\int_M
      \bigl(|\nabla^2H|_g^2+|\nabla H|_g^4\bigr)\,dV_g.
\]
Under $(g,H)\mapsto(\lambda^2g,\lambda^2H)$, both
$\|H\|_{L^\infty(g)}$ and $\mathfrak D_g(H)$ remain unchanged.
In particular no pointwise bound on the derivatives of $H$ is imposed.

\paragraph{Lemma 1 (uniform finite-variation remainders).}
Suppose $g$ is a smooth metric on a closed four-manifold and
\[
 \int_M|\operatorname{Rm}_g|^2\,dV_g\le KA.
\]
If
\[
 \|H\|_{L^\infty(g)}\le1,\qquad \mathfrak D_g(H)\le A,
\]
then $g_t=g+tH$ is positive definite for $|t|\le t_0$, where $t_0>0$
is dimensional, and
\begin{equation}
 \left|\frac{d^m}{dt^m}P(g_t)\right|\le C_K A,
 \qquad m=2,3,\quad |t|\le t_0.
 \label{eq:o11-var-remainder-20260929}
\end{equation}
The constants do not require a diameter bound, injectivity-radius bound,
Sobolev inequality, or noncollapsing hypothesis.

\begin{proof}
The amplitude bound makes $g_t$ uniformly equivalent to $g$. In terms of
the background connection $\nabla=\nabla^g$, the difference tensor is
\[
 (\nabla^{g_t}-\nabla^g)^k{}_{ab}
 =\frac t2(g_t)^{kc}
   \bigl(\nabla_aH_{bc}+\nabla_bH_{ac}-\nabla_cH_{ab}\bigr).
\]
The curvature difference is the covariant derivative of this tensor plus
its quadratic product. Thus it consists of contractions of
\[
 t\,g_t^{-1}\nabla^2H
 \quad\text{and}\quad
 t^2g_t^{-1}g_t^{-1}\nabla H\nabla H.
\]
Changing contractions or lowering indices also involves the background
curvature and the uniformly equivalent metrics. Differentiating in $t$
through order three only inserts further uniformly bounded inverse
metrics and copies of $H$. Consequently the curvature, traceless Ricci
tensor, and their required $t$-derivatives are bounded pointwise by
\[
 C_m\bigl(|\operatorname{Rm}_g|+|\nabla^2H|+|\nabla H|^2\bigr).
\]
The volume density relative to $dV_g$, together with its first three
$t$-derivatives, is uniformly bounded. Differentiating
$|\operatorname{Ric}^{\circ}_{g_t}|_{g_t}^2dV_{g_t}$ and applying Young's
inequality to the cross terms bounds each integral in
\eqref{eq:o11-var-remainder-20260929} by
\[
 C\int_M\bigl(|\operatorname{Rm}_g|^2
                  +|\nabla^2H|^2+|\nabla H|^4\bigr)\,dV_g
 \le C_KA.
\]
No coordinate neighborhood of a uniform radius is used.
\end{proof}

\paragraph{Proposition 2 (weak stationarity and stability at scale $Q_j$).}
For every prescribed sequence $\varepsilon_j>0$ there are smooth metrics
$g_j$ on $M_j$ such that
\[
 s_{g_j}=-4,\qquad 0\le P(g_j)-p_j\le\varepsilon_j,
 \qquad P(g_j)-p_j\longrightarrow0,
 \qquad \mathcal E(g_j)-\alpha_j=o(h_j),
\]
and, uniformly over all smooth symmetric two-tensors satisfying
\[
 \|H\|_{L^\infty(g_j)}\le1,\qquad
 \mathfrak D_{g_j}(H)\le A_j,
\]
one has
\begin{equation}
 |DP_{g_j}[H]|=o(Q_j),\qquad
 D^2P_{g_j}[H,H]\ge-o(Q_j).
 \label{eq:o11-var-selection-20260929}
\end{equation}
Here $D^2P_{g_j}[H,H]$ means the second derivative along $g_j+tH$.
It is not being identified with an intrinsic Hessian at an assumed
critical point.

\begin{proof}
Choose any $\eta_j\downarrow0$ and set
\begin{equation}
 \delta_j=\min\left\{
       \varepsilon_j,j^{-1},\eta_j^3\frac{Q_j^3}{A_j^2}
                    \right\}.
 \label{eq:o11-var-delta-20260929}
\end{equation}
The definition of the infimum supplies a smooth metric with $P\le
p_j+\delta_j$. Replace it by the negative Yamabe representative and
scale to $s=-4$. The inherited conformal formula and scale invariance
preserve
\begin{equation}
             0\le P(g_j)-p_j\le\delta_j.
 \label{eq:o11-var-gap-20260929}
\end{equation}
Since $p_j\le8\pi^2Q_j$ and $\delta_j=o(Q_j)$, these metrics satisfy
$P(g_j)=O(Q_j)$. The two-sided Ricci estimate yields
\[
 R(g_j)=8\pi^2A_j+O(Q_j).
\]
At scalar curvature $-4$, $R=P+4\operatorname{Vol}$, and hence
\[
 \operatorname{Vol}(g_j)=2\pi^2A_j+O(Q_j).
\]
The curvature identity and Euler-characteristic estimate then give
\[
 \int_{M_j}|\operatorname{Rm}_{g_j}|^2\,dV_{g_j}=O(A_j).
\]
Thus Lemma 1 applies uniformly to the stated test class.

For such an $H$, put $F(t)=P(g_j+tH)$. Every $g_j+tH$ with
$|t|\le t_0$ is an admissible global competitor, so
\[
                  F(t)\ge p_j\ge F(0)-\delta_j.
\]
Taylor expansion for both signs of $t$, using $|F''|\le CA_j$, gives
\[
 |F'(0)|\le\frac{\delta_j}{t}+CA_jt,
 \qquad 0<t\le t_0.
\]
Choosing $t$ comparable to $\sqrt{\delta_j/A_j}$ gives
\begin{equation}
 |DP_{g_j}[H]|\le C\sqrt{A_j\delta_j}
       \le C\eta_j^{3/2}Q_j\sqrt{h_j}=o(Q_j).
 \label{eq:o11-var-first-20260929}
\end{equation}
To estimate the second variation, add the inequalities for $t$ and $-t$.
The bound $|F'''|\le CA_j$ gives
\[
 F''(0)\ge-\frac{2\delta_j}{t^2}-CA_jt.
\]
Taking $t$ comparable to $(\delta_j/A_j)^{1/3}$, which is at most a
constant times $\eta_jh_j$, yields
\begin{equation}
 D^2P_{g_j}[H,H]\ge-CA_j^{2/3}\delta_j^{1/3}
                        \ge-C\eta_jQ_j.
 \label{eq:o11-var-second-20260929}
\end{equation}
Because $\eta_j\to0$, this is the asserted lower bound $-o(Q_j)$.
Both estimates are uniform over the entire test class.

Finally, constant scalar curvature gives $\mathcal E=4P/R$. The Ricci
estimate implies
\[
 \mathcal E(g_j)=\frac{4P(g_j)}{8\pi^2A_j}+O(h_j^2).
\]
Comparing with the inherited formula for $\alpha_j$ and using
\eqref{eq:o11-var-gap-20260929} gives
\[
 0\le\mathcal E(g_j)-\alpha_j
       =O(h_j^2)+O(\delta_j/A_j)=o(h_j).
\]
The prescribed absolute accuracy and convergence to zero follow from
\eqref{eq:o11-var-delta-20260929}.
\end{proof}

The conclusion is a selection statement about sufficiently accurate
global near-minimizers. It does not say that every competitive sequence
with $P=O(Q)$ is approximately stationary or stable.

\paragraph{The resulting weak equation.}
Write $W_+(g)=\int|W_g^+|^2\,dV_g$ and $T_+=2\chi+3\tau$.
Chern--Weil gives
\[
                 P=4W_++S/12-8\pi^2T_+.
\]
At a closed constant-scalar-curvature metric with $s=-4$, the scalar
square variation is
\[
 DS_g[H]=8\int_M\langle\operatorname{Ric}^{\circ}_g,H\rangle_g\,dV_g.
\]
Indeed the general scalar variation, integrated against the constant
scalar curvature, has no surviving divergence terms; its gradient is
$-2s\operatorname{Ric}^{\circ}$.
Proposition 2 therefore yields
\begin{equation}
 4DW_{+,g_j}[H]
 +\frac23\int_{M_j}
       \langle\operatorname{Ric}^{\circ}_{g_j},H\rangle_{g_j}\,dV_{g_j}
       =o(Q_j)
 \label{eq:o11-var-weak-equation-20260929}
\end{equation}
uniformly in the critical test class. Keeping the Weyl first variation
in this form avoids importing a separate Bach-tensor sign or norm
convention.

\paragraph{Shrinking supports.}
For a test tensor of fixed amplitude on a region of radius $r$ with
controlled rescaled geometry, the estimates
\[
 |\nabla H|=O(r^{-1}),\qquad |\nabla^2H|=O(r^{-2}),\qquad
 \operatorname{Vol}(\operatorname{supp}H)=O(r^4)
\]
give $\mathfrak D_g(H)=O(1)$, even when $r\to0$.
This illustrates how a microscopic test can satisfy the budget. The
controlled rescaled geometry in this illustration is not a hypothesis
of Proposition 2. A fixed constant multiple of the displayed budget
can also be allowed by multiplying $H$ by a fixed sufficiently small
constant.

\paragraph{Corollary 3 (disjoint-region selection without smooth precompactness).}
Let $g_j$ be selected as in Proposition 2. For each $j$, let
$H_{ji}$ be finitely many smooth symmetric two-tensors with pairwise
disjoint supports, and suppose
\[
 \|H_{ji}\|_{L^\infty(g_j)}\le1,
 \qquad \sum_i\mathfrak D_{g_j}(H_{ji})\le A_j.
\]
Then
\begin{equation}
 \sum_i|DP_{g_j}[H_{ji}]|=o(Q_j),\qquad
 \sum_i\bigl(-D^2P_{g_j}[H_{ji},H_{ji}]\bigr)_+=o(Q_j).
 \label{eq:o11-var-multiregion-20260929}
\end{equation}

\begin{proof}
Choose signs $\epsilon_i\in\{-1,1\}$ to make
$DP[\epsilon_iH_{ji}]\ge0$, and put $H=\sum_i\epsilon_iH_{ji}$.
Disjointness makes the amplitude bound persist and the derivative
budgets add. The linearity of the first variation and Proposition 2 give
the first conclusion.

For the second conclusion, sum only the tensors whose second variations
are negative. The functional $P$ is the integral of a local expression
in the second jet of the metric. Mixed second variations therefore
vanish for disjoint supports. The second variation along the sum is
exactly the sum of the negative second variations. Proposition 2 again
gives the result.
\end{proof}

This corollary assumes neither a fixed smooth compact family of cell
metrics nor uniform collar geometry or cell scale. It does require the
actual test tensors to satisfy the total critical budget. In particular,
a family of defects admitting definite descent per unit defect cannot
have positive $Q_j$-density if its descent tests satisfy that budget.
The existence of such tests for all potentially inexpensive defects has
not been proved here.

\paragraph{An unresolved instability criterion for the sharp coefficient.}
The following is a sufficient criterion, not an established curvature
inequality. Suppose that for every fixed $\epsilon\in(0,1)$ there are
$c_\epsilon>0$ and $j_\epsilon$ such that every smooth metric on $M_j$,
$j\ge j_\epsilon$, satisfying
\[
 s_g=-4,\qquad P(g)\le8\pi^2(1-\epsilon)Q_j
\]
admits a smooth symmetric tensor with
\[
 \|H\|_{L^\infty(g)}\le1,\qquad \mathfrak D_g(H)\le A_j,
\]
for which at least one of
\begin{equation}
 |DP_g[H]|\ge c_\epsilon Q_j,
 \qquad
 D^2P_g[H,H]\le-c_\epsilon Q_j
 \label{eq:o11-var-open-instability-20260929}
\end{equation}
holds. Then
\[
 \frac{p_j}{8\pi^2Q_j}\longrightarrow1,
 \qquad \frac{\alpha_j}{h_j}\longrightarrow4.
\]

\begin{proof}[Proof of the conditional implication]
If the lower limit of $p_j/(8\pi^2Q_j)$ were less than one, choose a
subsequence with a fixed strict gap below one. The selected metrics of
Proposition 2 remain below a slightly weaker fixed threshold, because
$P(g_j)-p_j=o(Q_j)$. Both alternatives in
\eqref{eq:o11-var-open-instability-20260929} contradict the corresponding
uniform estimates in Proposition 2. Hence the lower limit is at least
one. The recovery upper bound gives the first limit, and the inherited
variational comparison gives the second.
\end{proof}

Proving \eqref{eq:o11-var-open-instability-20260929}, or constructing
genuine metrics with a smaller first-order cost, remains the substantive
unrestricted problem. In particular, one cannot insert
$H=\operatorname{Ric}^{\circ}$ in the preceding weak equation without
checking its amplitude and critical derivative budget. Nor do these
selection statements establish smooth convergence, a bubble-tree
decomposition, or the signed-curvature compensation inequality.

\end{document}
